\chapter{The verifier}
\newcommand{\cEightCode}[1]{\texttt{\scriptsize #1}}

\section*{Why this matters for us}
The verifier is the part of Shelfie that decides whether money moves. Everything the pitch promises (``paid only when the photo is true'') rests on it. It is also the part judges will attack with questions about fraud. Understanding it deeply, including what it \emph{cannot} prove, lets any teammate answer those questions calmly. The details come from \texttt{diseno/verifier-rules.md}. As everywhere in this guide, nothing is built yet.

\section{Design principle: a deterministic gate around a probabilistic reader}
Two kinds of components appear in the verifier:
\begin{itemize}
  \item A \term{vision model} (an AI that looks at images) is \emph{probabilistic}: it gives an answer that is usually right, sometimes wrong, and can differ between two runs.
  \item A \term{deterministic} rule is ordinary code: the same input always gives the same output, and you can explain why.
\end{itemize}
Shelfie's rule: \textbf{the model only reads; plain code decides.} The model reports what it sees (a code, a price, a product name). Then rules compare those readings with what we expect. Payment depends on rules, never on the model's opinion.

\begin{analogybox}{Court clerk and judge}
The vision model is a clerk who reads out what is written on a document. The gates are the judge who applies the law to those words. If the clerk misreads a digit, the judge's explicit rules (does the code match? is the price plausible?) still catch most mistakes, and every decision can be explained afterwards.
\end{analogybox}

Another design rule: \textbf{the model never sees the expected code or the expected price.} If it knew the answer, it might ``confirm'' what we hope to see (confirmation bias). It reports blindly and our code compares.

\section{The challenge code}
\subsection{What problem it solves}
A worker is paid for a photo. The cheapest cheat is to reuse a photo: an old one, a stock image, a screenshot, or a photo taken at home. How can the verifier know that the photo was taken \emph{now} and \emph{for this task}? Ask the worker to include something that did not exist before the task started. That is the \term{challenge code}.

\begin{conceptbox}{Challenge code}
A random, single-use, short-lived 6-character code, issued when the worker claims a task. The worker must show it in the photo, next to the price tag. It proves \emph{freshness} (the photo is newer than the claim) and \emph{task binding} (the photo belongs to this task). It does \textbf{not} prove location.
\end{conceptbox}

\subsection{Format and alphabet}
\begin{itemize}
  \item 6 characters, displayed as \code{XXX-XXX}, for example \code{K7Q-4XM}, in a big monospace font.
  \item Alphabet: \code{ABCDEFGHJKMNPQRSTUVWXYZ23456789}. The characters \textbf{0, O, 1, I, L} are left out because they are easily confused in handwriting and photos (zero or letter O? one or letter L?). That leaves 31 characters.
  \item $31^6 = 887{,}503{,}681$ combinations, so guessing is hopeless (the design files say ``about $30^6$'', which is the same order). Codes come from a cryptographically secure random generator, not from the task id, so they cannot be predicted.
  \item Time to live: 300 s by default. An expired claim returns the task to OPEN.
  \item A new claim attempt always gets a \emph{new} code. A rejected photo burns its code.
\end{itemize}

\subsection{Why the phone cannot show its own screen}
The phone that takes the photo is busy being a camera. It cannot display the code and photograph itself at the same time. So the code has to be \emph{physically} next to the price tag. Two modes exist, chosen in the pack config:
\begin{itemize}
  \item \code{handwritten} (default): the worker reads the code on the phone, writes it on a piece of paper and holds it next to the tag.
  \item \code{card-pool}: printed cards with pre-generated codes; the engine hands out a card number. We use this for the stage demo for reliability. In the card-pool mode we print 20 large cards with high contrast.
\end{itemize}
A third mode, \code{none}, exists for Plan B where agents, not humans, do the work.

\begin{figure}[htbp]
\centering
\begin{tikzpicture}
% shelf
\fill[cGrey!15] (0,0) rectangle (13,0.25);
\fill[cGrey!25] (0,2.9) rectangle (13,3.1);
% price tag
\node[draw=cBlue,thick,fill=white,rounded corners=2pt,minimum width=3.6cm,minimum height=1.7cm,align=center,font=\sffamily] (tag) at (3.2,1.7) {\small Example Pale Ale\\\small 0,33 l\\[2pt]\Large\bfseries 2,49 EUR};
% bottle
\fill[cOrange!60] (0.8,0.25) rectangle (1.5,2.1);
\fill[cOrange!60] (1.0,2.1) rectangle (1.3,2.7);
% paper with code
\node[draw=cGrey,thick,fill=white,minimum width=3.4cm,minimum height=1.7cm,rotate=-4,align=center] (paper) at (7.6,1.6) {\ttfamily\bfseries\LARGE K7Q-4XM};
\node[font=\sffamily\scriptsize,text=cGrey] at (7.6,0.5) {paper or printed card};
% hand
\fill[cOrange!25,rounded corners=4pt] (8.9,0.9) rectangle (9.7,1.5);
% label arrows
\node[font=\sffamily\scriptsize,text=cBlue,align=center] at (3.2,3.5) {price tag\\(read by the model)};
\node[font=\sffamily\scriptsize,text=cGreen!60!black,align=center] at (7.6,3.5) {challenge code\\(compared by the rules)};
% camera frame
\draw[thick,cPurple,dashed,rounded corners=4pt] (0.4,-0.35) rectangle (12.6,3.95);
\node[font=\sffamily\scriptsize,text=cPurple,anchor=south east] at (12.6,-0.35) {one photo: tag and code both in frame};
\end{tikzpicture}
\caption{What a valid photo contains (schematic, not a real photo). The code is not shown on the phone; it is on paper or a card in the frame.}
\end{figure}

\begin{warningbox}{Freshness is not location}
The code proves the photo is new and for this task. A worker could still photograph the same product in a different shop. We state this limit in the pitch (ADR-12). Roadmap: store-owned QR codes or geofenced tasks.
\end{warningbox}

\section{The verification pipeline and gates G1 to G7}
The verifier receives a request (the photo plus the task context: expected code, product, time window) and produces one \term{verdict}. The verdict says \textbf{pass only if every gate passes}. There is no weighted score; scores exist only inside a gate, against that gate's own threshold. All gates are evaluated and reported, so the worker and the dashboard see every failure at once. The one exception: if G1 fails, the (costly) vision call is skipped.

\begin{figure}[htbp]
\centering
\resizebox{\textwidth}{!}{%
\begin{tikzpicture}
\node[service,minimum width=2.3cm] (in) at (0,0) {Request:\\photo + task};
\node[bad,fill=cOrange!10,draw=cOrange,minimum width=1.8cm] (g1) at (3,0) {G1\\image\_ok};
\node[agent,minimum width=2.1cm] (vis) at (6,0) {Vision model\\one call};
\node[actor,minimum width=1.8cm] (g2) at (9,1.1) {G2\\code\_match};
\node[actor,minimum width=1.8cm] (g3) at (11.4,1.1) {G3\\readable};
\node[actor,minimum width=1.8cm] (g4) at (13.8,1.1) {G4\\product\_match};
\node[actor,minimum width=1.8cm] (g5) at (9,-1.1) {G5\\price\_sanity};
\node[actor,minimum width=1.8cm] (g6) at (11.4,-1.1) {G6\\timing};
\node[actor,minimum width=1.8cm] (g7) at (13.8,-1.1) {G7\\not\_duplicate};
\node[good,minimum width=2.4cm] (out) at (16.8,0) {Verdict:\\PASS only if\\all gates pass};
\draw[flow] (in)--(g1); \draw[flow] (g1)--(vis);
\draw[flow] (vis.east) -- ++(0.5,0) |- (g2.west);
\draw[flow] (g2)--(g3); \draw[flow] (g3)--(g4);
\draw[flow] (g4.east) -- ++(0.35,0) |- (out.north west);
\draw[flow] (vis.east) ++(0.5,0) |- (g5.west);
\draw[flow] (g5)--(g6); \draw[flow] (g6)--(g7);
\draw[flow] (g7.east) -- ++(0.35,0) |- (out.south west);
\draw[flow,draw=cRed,dashed] (g1.north) -- ++(0,1.3) -| node[lbl,pos=0.15]{G1 fails: skip vision call} (out.north);
\node[font=\sffamily\scriptsize,text=cGrey,anchor=north] at (11.4,-2.1) {G6 and G7 need server-side data: engine receive time and earlier image hashes};
\end{tikzpicture}}
\caption{The verification pipeline. The vision call is the only impure step; the seven rules are pure functions of the stored extraction, so they can be tested with fixed JSON and no model.}
\end{figure}

\subsection{The seven gates, each with a pass and a fail example}
Use \code{Example Pale Ale 0.33l}, expected code \code{K7Q-4XM}, as the running example.

{\small
\renewcommand{\arraystretch}{1.15}
\begin{longtable}{@{}p{2.3cm}p{3.5cm}p{4.2cm}p{4.2cm}@{}}
\toprule
\textbf{Gate} & \textbf{What it checks} & \textbf{Pass example} & \textbf{Fail example (reason code)}\\
\midrule
\endhead
\textbf{G1} \cEightCode{image\_ok} & Decodable JPEG or PNG; at least 800$\times$600 pixels; not too big; not a near-blank image. & 1600$\times$1200 JPEG, 400 KB. & 640$\times$480 image (\cEightCode{image\_too\_small}).\\
\textbf{G2} \cEightCode{code\_match} & Code read from the photo equals the expected code (edit distance at most 0) and read confidence at least 0.7. & Reads \cEightCode{K7Q-4XM}, confidence 0.93. & Reads \cEightCode{K7Q-4XN} (\cEightCode{code\_mismatch}); no code visible (\cEightCode{code\_missing}).\\
\textbf{G3} \cEightCode{readable} & Price parses to a number (by our parser); currency EUR or absent; product text not empty; overall confidence at least 0.6. & Text \cEightCode{2,49 EUR} parses to 2.49. & Glare hides the price (\cEightCode{price\_unreadable}).\\
\textbf{G4} \cEightCode{product\_match} & Valid EAN-13 barcode equals the task EAN, or name similarity to name or an alias at least 0.8. & Reads \cEightCode{Example Pale Ale 0,33 l}. & Reads another beer (\cEightCode{product\_mismatch}).\\
\textbf{G5} \cEightCode{price\_sanity} & Between 0.20 and 60.00 EUR; within 60\% of a reference price if the job gives one; at most 2 decimals. & 2.49 EUR. & 249.00 EUR (\cEightCode{price\_out\_of\_range}).\\
\textbf{G6} \cEightCode{timing} & Engine's receive time is at least 5 s after the claim and before the claim expires and before the task window ends. & Received 40 s after the claim. & Received after the 300 s TTL (\cEightCode{claim\_expired}).\\
\textbf{G7} \cEightCode{not\_duplicate} & Image hash not seen in the same job; perceptual hash not too close to earlier accepted photos. & A new photo. & The same file submitted twice (\cEightCode{duplicate\_image}).\\
\bottomrule
\end{longtable}
}

A few details that often confuse people:
\begin{itemize}
  \item \textbf{Edit distance} counts how many characters you must change to turn one string into another. \code{K7Q-4XM} to \code{K7Q-4XN} is distance 1. The default allows 0, because our alphabet has no confusable characters.
  \item \textbf{G6 uses the engine's clock}, never the phone's clock and never the photo's built-in EXIF date, because both are editable. EXIF time is only logged as advisory.
  \item A \textbf{perceptual hash} is a small fingerprint of what an image \emph{looks like}, so that a slightly cropped or recompressed copy of a photo is still recognised as the same.
  \item The worker sees a plain-language message per failure, for example ``We could not read the code K7Q-4XM in your photo. Hold the paper closer and try again.''
\end{itemize}

\begin{examplebox}{Cheat attempts we rehearse (all must be rejected)}
\begin{enumerate}
  \item An old camera-roll photo of the tag: no code in frame, G2 \code{code\_missing}.
  \item The right tag with a wrong code: G2 \code{code\_mismatch}.
  \item A photo of a different product with the right code: G4 \code{product\_mismatch}.
  \item The price covered or blurred: G3 \code{price\_unreadable}.
  \item The same photo submitted twice: G7 \code{duplicate\_image}.
\end{enumerate}
In every case the check was still paid to the verifier (1 tADA); the worker is paid nothing.
\end{examplebox}

\section{The vision model contract}
One call per submission. The prompt lives in the domain pack. The model must return \emph{only} JSON in a fixed shape (illustrative, from the design file):

\begin{lstlisting}
{ "code":    { "text": "K7Q-4XM", "confidence": 0.93 },
  "price":   { "text": "2,49 EUR", "value": 2.49,
               "currency": "EUR", "confidence": 0.9 },
  "product": { "text": "Example Pale Ale 0,33 l", "confidence": 0.88 },
  "barcode": { "digits": "4000000000000", "confidence": 0.4 },
  "scene":   { "isShelfOrPriceTag": true, "peopleVisible": false,
               "glare": "low", "blur": "none" },
  "advisory":{ "tamperSuspected": false, "looksLikeScreenPhoto": false },
  "overallConfidence": 0.88 }
\end{lstlisting}

Rules for the adapter, in plain words:
\begin{enumerate}
  \item Unreadable fields are \code{null} with confidence 0, never guessed.
  \item The JSON is parsed strictly; on a parse failure the call is retried once with a repair instruction, then the verifier answers ``503 vision unavailable'' (no verdict, and if possible no charge).
  \item German number formats such as \code{2,49} are parsed by \emph{our} parser. If the model's \code{value} and our parse disagree, ours wins.
  \item Images are downscaled to at most 1600 pixels on the long side for cost and speed.
  \item Which model and provider is used is \emph{configuration} (ADR-07), chosen by the 20-photo test below, not by preference. Whether the shell's own language model can read images is unknown (verify Friday).
\end{enumerate}

\section{Thresholds}
Every number is in \code{pack.config.json}, so tuning needs no code change.

\begin{center}
\small
\begin{tabularx}{\textwidth}{@{}p{5.6cm}p{1.9cm}X@{}}
\toprule
\textbf{Parameter} & \textbf{Start} & \textbf{How it is tuned}\\
\midrule
\cEightCode{code\_match.maxEditDistance} & 0 & Raise to 1 only if handwritten codes fail but wrong-code attacks still fail.\\
\cEightCode{code\_match.minConfidence} & 0.7 & Lower if too many false rejects, never below 0.5.\\
\cEightCode{readable.minConfidence} & 0.6 & Same.\\
\cEightCode{product\_match.minTokenScore} & 0.8 & Add aliases first; lower the threshold last.\\
\cEightCode{price\_sanity.minEur / maxEur} & 0.20 / 60 & Set per demo product.\\
\cEightCode{price\_sanity.maxDeviationFromRef} & 0.6 & Only when the job gives a reference price.\\
\cEightCode{timing.minSecondsAfterClaim} & 5 & Fixed.\\
\cEightCode{image\_ok.minWidth / minHeight} & 800 / 600 & Fixed.\\
Verify latency budget & p95 under 8 s & Faster model if exceeded (excludes the x402 step).\\
\bottomrule
\end{tabularx}
\end{center}

\begin{conceptbox}{False accept versus false reject}
A \term{false accept} is a bad photo that passes (we pay for fraud). A \term{false reject} is a good photo that fails (an honest worker is annoyed). For a payment gate, false accepts are the safety failure. That is why the acceptance targets below demand \emph{zero} false accepts but tolerate one false reject.
\end{conceptbox}

\section{The 20-photo evaluation}
Goal: choose the model and thresholds with evidence and obtain an honest number for the pitch. Owner C runs it on Friday as soon as the verifier answers (target: before milestone M2). Capture takes about 45 minutes with two people: printed price tags for a real product on our own shelf (or a fridge) plus a few from a shop with permission; \textbf{no people in frame}; each photo gets a throwaway code on paper.

\begin{center}
\small
\begin{tabularx}{\textwidth}{@{}p{1.1cm}p{6.6cm}X@{}}
\toprule
\textbf{Count} & \textbf{Case} & \textbf{Expected}\\
\midrule
6 & Valid: correct code, readable tag, right product, different angles and light & PASS\\
2 & Valid but hard: glare or sharp angle & PASS (a reject here is a false reject)\\
3 & Wrong code (different valid format; one off by one character) & FAIL \code{code\_mismatch}\\
2 & No code in frame (camera-roll style) & FAIL \code{code\_missing}\\
2 & Wrong product, right code & FAIL \code{product\_mismatch}\\
2 & Tag blurred or covered & FAIL unreadable / low confidence\\
1 & Absurd edited price (e.g.\ 249.00) & FAIL \code{price\_out\_of\_range}\\
1 & Same image as a valid one, resubmitted & FAIL \code{duplicate\_image}\\
1 & Photo of a screen showing a tag & Record what happens (not required to detect)\\
\bottomrule
\end{tabularx}
\end{center}

\begin{figure}[htbp]
\centering
\begin{tikzpicture}
\def\u{0.72}
\fill[cGreen!70] (0,0) rectangle (6*\u,0.8);
\fill[cGreen!40] (6*\u,0) rectangle (8*\u,0.8);
\fill[cRed!60] (8*\u,0) rectangle (11*\u,0.8);
\fill[cRed!45] (11*\u,0) rectangle (13*\u,0.8);
\fill[cRed!35] (13*\u,0) rectangle (15*\u,0.8);
\fill[cOrange!50] (15*\u,0) rectangle (17*\u,0.8);
\fill[cOrange!70] (17*\u,0) rectangle (18*\u,0.8);
\fill[cPurple!50] (18*\u,0) rectangle (19*\u,0.8);
\fill[cGrey!40] (19*\u,0) rectangle (20*\u,0.8);
\foreach \a/\b/\t in {0/6/6 valid,6/8/2 hard,8/11/3 wrong code,11/13/2 no code,13/15/2 wrong prod.,15/17/2 blurred,17/18/1,18/19/1,19/20/1}
  \node[font=\sffamily\tiny] at ({(\a+\b)/2*\u},0.4) {\t};
\draw[decorate,decoration={brace,amplitude=4pt}] (0,0.95) -- (8*\u,0.95) node[midway,above=3pt,font=\sffamily\scriptsize] {8 should pass};
\draw[decorate,decoration={brace,amplitude=4pt}] (8*\u,0.95) -- (20*\u,0.95) node[midway,above=3pt,font=\sffamily\scriptsize] {12 should fail (screen photo included)};
\end{tikzpicture}
\caption{Composition of the 20-photo test set. The last three blocks: absurd price, duplicate, screen photo.}
\end{figure}

\subsection{Acceptance targets}
\begin{center}
\small
\begin{tabularx}{\textwidth}{@{}p{5.2cm}p{3cm}X@{}}
\toprule
\textbf{Metric} & \textbf{Target} & \textbf{If missed}\\
\midrule
False accepts on the 12 should-fail cases & \textbf{0 of 12} & Tighten thresholds or change model. Blocking for payments.\\
False rejects on the 8 valid cases & at most 1 of 8 & Add aliases, loosen confidence, improve prompt, try another model.\\
Price value exact on valid cases & at least 7 of 8 & Prompt or model change.\\
Code read exact on valid cases & at least 7 of 8 & Switch to \code{card-pool} mode with large printed cards.\\
Latency p95 & under 8 s & Faster model, smaller image.\\
Cost per check & recorded & A number for the pitch.\\
\bottomrule
\end{tabularx}
\end{center}
A runner script posts every photo to a dry (no-payment) verify endpoint, compares the results with labels, prints a confusion table and stores the raw extractions, so thresholds can be re-tuned offline without new model calls. With only 20 photos, the estimate is rough. The pitch says ``tested on 20 photos'', \emph{never} ``99\% accurate''.

\section{Manual override: a stage safety net, clearly flagged}
If the vision provider misreads on stage, a human operator may keep the demo flowing. This is a rescue tool, not a product feature, and it is limited by strict rules:

\begin{itemize}
  \item \textbf{Off by default} (\code{ALLOW\_DEMO\_OVERRIDE=false}). It must be off for every evidence run (milestones M1 to M5) and for recorded video takes. Turning it on is itself written to the ledger.
  \item \textbf{Admin only:} one endpoint behind \code{X-Admin-Token}; there is no button on the public worker page.
  \item Every override creates a verdict with source \code{MANUAL\_OVERRIDE}, a \textbf{required reason}, and a ledger entry flagged \code{demoOverride}.
  \item \textbf{Visible everywhere:} the dashboard row shows a red ``DEMO OVERRIDE'' badge, the report footnotes the row, and the Devfolio text says the override exists and which rows (if any) used it.
  \item It can \emph{never} rescue a photo that failed G1 or G7, so one photo cannot be paid twice. It may only turn a FAIL into a PASS for the first accepted attempt of a task, one override per submission.
\end{itemize}

\begin{warningbox}{Credibility depends on never using it silently}
The whole claim is ``paid only when the photo is true''. An unlabelled override would make the claim false. If the override was needed in an evidence run, re-record instead.
\end{warningbox}

\section{Privacy}
Stores are semi-public places with staff and customers. Rules:
\begin{itemize}
  \item The task tells workers: \textbf{shelf and tag only, no people}.
  \item The vision model reports \code{peopleVisible}. If \code{true}, an advisory note is attached to the verdict and the operator sees a warning. It is an \emph{advisory flag}: it does \textbf{not} block the payment today. Automatic rejection or cropping is roadmap. (An earlier answer draft claimed faces were cropped; the review report corrected this.)
  \item Images are stored only on the hub laptop during the event (for the video and disputes) and deleted afterwards.
  \item Workers' data: a random worker id and a wallet address, nothing else.
  \item Open question: whether stores permit photographing; we ask permission for the video or use our own shelf.
\end{itemize}

\section{What the verifier cannot prove}
Say these limits yourself before a judge does:
\begin{center}
\small
\begin{tabularx}{\textwidth}{@{}p{4.6cm}X@{}}
\toprule
\textbf{Not proven} & \textbf{Why / what we do}\\
\midrule
Location (``the photo was taken in this store'') & A worker can shoot the same product elsewhere. Roadmap: store-owned QR codes, geofenced tasks.\\
The tag sits on the right shelf slot & Accepted risk for the hackathon.\\
Authenticity beyond freshness and readability & A live-made fake is possible in theory; the claim is that cheating costs more than the payout.\\
Identity of the worker & Only a wallet address; no reputation or sybil protection yet.\\
The model never errs & It can misread; G2 to G5 catch many errors, but 20 test photos give only a rough rate.\\
Screenshots and photos of screens & Not reliably detected by the gates. The code is generated at claim time, which makes a prepared fake hard; \code{looksLikeScreenPhoto} is advisory only.\\
\bottomrule
\end{tabularx}
\end{center}

\begin{keypoints}
\begin{itemize}
  \item The verifier is a \textbf{deterministic gate around a probabilistic reader}: the model reads, code decides; the model never sees the expected code or price.
  \item The \textbf{challenge code}: 6 characters \code{XXX-XXX}, alphabet without 0/O/1/I/L ($31^6$ combinations), single use, 300 s, shown on the phone but written on paper or a printed card.
  \item \textbf{G1 to G7}: image\_ok, code\_match, readable, product\_match, price\_sanity, timing, not\_duplicate. Pass only if all seven pass.
  \item Pass or fail, each check costs 1 tADA, paid to the verifier via x402.
  \item The 20-photo test must show \textbf{0 false accepts} of 12 bad cases; say ``tested on 20 photos''.
  \item The manual override is admin-only, red-badged, off in evidence runs, and cannot rescue G1 or G7 failures.
  \item Limits: freshness not location; \code{peopleVisible} is advisory; no identity checks.
\end{itemize}
\end{keypoints}

\begin{quizbox}
\begin{enumerate}
  \item Why is the code alphabet missing the characters 0, O, 1, I and L?
  \item Why can the phone that takes the photo not show the code on its own screen?
  \item Which gate fails for an old camera-roll photo, and with which reason code?
  \item Why does G6 use the engine's clock rather than the photo's date?
  \item What are the acceptance targets for false accepts and false rejects on the 20 photos?
  \item Name three rules that limit the manual override.
\end{enumerate}
{\small\color{cGrey}\textbf{Answers.}
(1) They are easy to confuse in handwriting and photos.
(2) The phone is the camera and cannot photograph itself; the code must be physically in the frame on paper or a card.
(3) G2 \code{code\_match}, reason \code{code\_missing}.
(4) EXIF and device clocks can be edited; the server's receive time cannot.
(5) 0 of 12 false accepts; at most 1 of 8 false rejects.
(6) Any three of: off by default; admin token only; mandatory reason; red badge and ledger flag; never in evidence runs; cannot override G1 or G7; one per submission.}
\end{quizbox}
